% SOURCE_AUTHORITY: sga5_fr_workpass.tex, lines 7941--8004 (LF-normalized unit).
% SOURCE_SHA256: 791F4EFFC5E02832D5D77ED03518C8156D6F07E4C8238B03545DB93D883FBB28
\subsubsection*{A.2. Caso de los sistemas proyectivos}
Sea ahora \(X=(X_n)_{n\in\mathbb Z}\) un sistema proyectivo de objetos de una categoría abeliana \(\mathcal C\), tal que
\[
X_i=0\qquad(i<0).
\]
Si \(T\) denota un \(\delta\)-funtor exacto de \(\mathcal C\) en una categoría abeliana \(\mathcal D\), \(T\) se prolonga de manera natural a un \(\delta\)-funtor exacto, denotado todavía por \(T\):
\[
\uHom(\mathbb Z^\circ,\mathcal C)\longrightarrow \uHom(\mathbb Z^\circ,\mathcal D).
\]
Por tanto, dotando a \(X\) de su filtración natural de sistema proyectivo (EGA \(0_{\mathrm{III}}\) 13.4.1), se obtiene una sucesión espectral:
\[
\mathbb{E}(X)=\mathbb{E}(X_n)_{n\in\mathbb Z},
\]
donde los \(\mathbb{E}(X_n)\) denotan las sucesiones espectrales asociadas a \(T\) y a los \(X_n\) dotados de la filtración inducida.

\paragraph{Propiedades de \(\mathbb{E}(X)\).}
a) Para todo entero \(n\), la filtración del término límite \(E^n(X)\) en la sucesión espectral \(\mathbb{E}(X)\) coincide con su filtración natural de sistema proyectivo cuando \(X\) es estricto.

En efecto, denotando por la letra \(F\) (como en A.1) la primera filtración y por la letra \(\Phi\) la segunda, se tiene, para todo entero \(\alpha\):
\[
F^pE^n(X)_\alpha
=\Ker\bigl(T^n(X_\alpha)\longrightarrow T^n(X_\alpha/F^pX_\alpha)\bigr)
\]
\[
=
\begin{cases}
0,& \alpha\le p-1,\\
\Ker\bigl(T^n(X_\alpha)\longrightarrow T^n(X_{p-1})\bigr),& \alpha\ge p-1.
\end{cases}
\]
En ambos casos se reconoce la expresión de \(\Phi^pE^n(X)\).

Esto muestra, en particular, que, para \(\beta\ge\alpha\ge p\), el morfismo de transición
\[
E_\infty^{pq}(X_\beta)\longrightarrow E_\infty^{pq}(X_\alpha)
\]
que se identifica con
\[
\gr^pT^{p+q}(X_\beta)\longrightarrow
\gr^pT^{p+q}(X_\alpha),
\]
es un monomorfismo, por EGA \(0_{\mathrm{III}}\) 13.4.1.3.

b) Fijado \(\alpha\in\mathbb Z\), es claro (cf. la figura adjunta) que, puesto que \(F^i(X_\alpha)=0\) para \(i\ge\alpha+1\), se tiene:
\[
\begin{tikzpicture}[baseline=(current bounding box.center),x=1.0cm,y=1.0cm,>=latex]
\draw (-1.27,0) -- (3.81,0);
\draw (0,-0.5) -- (0,2.77);
\draw (1.78,-1.05) -- (1.78,2.5);
\node[below left] at (0,0) {$0$};
\node[below right] at (1.78,0) {$\alpha$};
\draw[->] (1.38,2.22) -- (1.38,1.33);
\draw[->] (1.38,1.33) -- (2.68,1.33);
\draw (1.38,2.22) -- (2.68,1.33);
\node[above] at (1.33,2.22) {$(p,q)$};
\node[left] at (1.38,1.78) {$1-r$};
\node[above] at (2.08,1.33) {$r$};
\end{tikzpicture}
\]
\[
B_r^{pq}(X_\alpha)=B_{r+1}^{pq}(X_\alpha)\quad\text{para }r>p,
\]
\[
Z_r^{pq}(X_\alpha)=Z_{r+1}^{pq}(X_\alpha)\quad\text{para }r>\alpha-p.
\]
